The Large Language models (LLMs) trained from massive text datasets allow them to be equipped with sufficient knowledge of different research areas. Recent studies have demonstrated that LLMs can assist vehicle agents in solving complex sequential decision-making tasks such as navigation. However, the performance on navigation tasks varies with different reasoning and prompting methods (e.g., Chain of Thought, Buffer of Thought, etc.) of the LLM agent. In this study, we compare 4 different LLM reasoning methods (i.e., ReAct, Chain of Thought, Tree of Thought, Buffer of Thought) with the Reactive baseline. The MiniGrid environment is selected to test the effectiveness of the 5 methods. More specifically, the average success rate, exploration steps, and total tokens used are applied as the performance metrics. Evaluation was done on three of the MiniGrid environments, including {\fontfamily{qcr}\selectfont Empty8$\times$8}, {\fontfamily{qcr}\selectfont LavaGapS6} and {\fontfamily{qcr}\selectfont LavaCrossingS9N2} while considering the difficulty levels. The experimental results demonstrate that most of the reasoning-enhanced LLM agents outperform the reactive baseline in success rate and exploration steps, with a certain compute trade-off.